\documentclass[11pt]{article}
\usepackage[a4paper,margin=27mm]{geometry}
\usepackage{amsmath,amssymb,mathtools}
\usepackage{fontspec,unicode-math}
\setmainfont{Latin Modern Roman}
\setmathfont{Latin Modern Math}
\usepackage{xurl}
\usepackage[hidelinks]{hyperref}
\newcommand{\sref}[2]{\hyperlink{source:#1:#2}{[S#2]}}
\newcommand{\cataloguescope}[1]{\paragraph{Mathematical question.}#1}
\hypersetup{pdftitle={UnsolvedMath Problem 9400137}}
\begin{document}
% UnsolvedMath ID 9400137; source catalogue code AMR-093-0137
\setcounter{section}{9400136}
\section{Does the converse of Wolstenholme's theorem hold for all natural numbers}\label{9400137}
{\small\sffamily UnsolvedMath ID 9400137 · Number Theory}
\subsection{Definitions and mathematical statement}

\paragraph{Shared notation.}$\mathbb N=\{1,2,\ldots\}$, and $\mathbb Z,\mathbb Q,\mathbb R,\mathbb C$ denote the integers, rationals, reals and complex numbers. Any other conventions are specified locally.


Wolstenholme's theorem states that every prime $p\ge5$ satisfies $\binom{2p-1}{p-1}\equiv1\pmod{p^3}$, where $\binom ab=a!/(b!(a-b)!)$. Its proposed converse has the same domain of integers at least five.
\paragraph{Statement recorded in the supplied source.}
For every integer $n\ge5$, does $\binom{2n-1}{n-1}\equiv1\pmod{n^3}$ imply that $n$ is prime? Equivalently, does every composite integer $n\ge5$ fail this congruence? \sref{9400137}{2}
\subsection{Short English statement}
Does the cubic-modulus binomial congruence characterize primes in the range of Wolstenholme’s theorem?
\subsection{Sources}
\begin{description}
\item[\hypertarget{source:9400137:1}{S1}] ULAM.AI and the UnsolvedMath Contributors, UnsolvedMath: A Curated Collection of Open Mathematics Problems, record 9400137 (AMR-093-0137). CC BY 4.0. \url{https://huggingface.co/datasets/ulamai/UnsolvedMath}.
\item[\hypertarget{source:9400137:2}{S2}] Jonathan Sondow, \emph{Extending Babbage\textquotesingle s (non-)primality tests}, arXiv:1812.07650 (2018) Equation (10) and converse discussion, pp. 4--5. \url{https://arxiv.org/pdf/1812.07650}.
\item[\hypertarget{source:9400137:3}{S3}] Wikipedia contributors, \emph{List of unsolved problems in mathematics}, revision 1366636767 (2026) Number theory, Does the converse of Wolstenholme's theorem hold for all natural numbers. \url{https://en.wikipedia.org/w/index.php?title=List_of_unsolved_problems_in_mathematics&oldid=1366636767#Number_theory}.
\end{description}
\label{end:9400137}
\end{document}
